\begin{ledgroupsized}[r]{120mm}
\footnotesize
\pstart
\noindent\textbf{Überlieferung:}
\pend
\end{ledgroupsized}
\begin{ledgroupsized}[r]{114mm}
\footnotesize
\pstart
\parindent -6mm
\makebox[6mm][l]{\textit{L}}%
Auszüge mit Bemerkungen aus
\indexN{\textls{Fermat}, Pierre de 1601\textendash1665}%
\textsc{P.\;de\;Fermat},
\cite{00533}\textit{Varia opera mathematica}, Toulouse 1679:
LH~XXXV~14,\;2 Bl.\;22--23. 
Ein Bogen 2\textsuperscript{o};
{Wasserzeichen} in Bl.\;22 (Papier aus Rom); %
oberer und unterer Rand beschnitten.
Dreieindrittel Seiten;
die oberen zwei Drittel von Bl.\;%
\href{\Scan{LH_035_14_02_022r}}
{22~r\textsuperscript{o}} (wie auch die benachbarten Bogen LH~XXXV~14,\;2 Bl.\;20--21 und 24--25) überliefern Auszüge mathematischen Inhalts aus demselben \cite{00533}Werk Fermats\indexN{\textls{Fermat}, Pierre de 1601\textendash1665}; sie werden in einem künftigen Band von \textit{LSB}~VII ediert.
Am oberen Rand von Bl.\;%
\href{\Scan{LH_035_14_02_022r}}
{22~r\textsuperscript{o}} findet sich folgender Vermerk von Leibnizens Hand:
\textit{Altera scheda meorum excerptorum ex Posthumis Fermatii}.
Die Diagramme im Text sind nicht der Vorlage entnommen.
\pend
\end{ledgroupsized}
\vspace{5mm}
\begin{ledgroup}
\footnotesize
\pstart
\noindent
\textbf{Datierungsgründe:}
Auf Grund des Wasserzeichens in Bl.\;22 kann eine Abfassung während der Aufenthalte in Rom (zweite Hälfte April bzw. Mitte Mai bis November 1689) angenommen werden.
\pend
\end{ledgroup}
\selectlanguage{latin}		%
\vspace{8mm}
\pstart
\noindent
\href{\Scan{LH_035_14_02_022r}}%
{\lbrack22~r\textsuperscript{o}\rbrack} 
Sequuntur Epistolae Fermatii\indexN{\textls{Fermat}, Pierre de 1601\textendash1665},
\edtext{scribit Mersenno,\indexN{\textls{Mersenne} (Mersennus), Marin 1588\textendash1648}}{%
\lemma{scribit Mersenno}%
\Cfootnote{%
\cite{02580}\indexN{\textls{Fermat}, Pierre de 1601\textendash1665}\textsc{P. Fermat}, Brief an \indexN{\textls{Mersenne} (Mersennus), Marin 1588\textendash1648}M. Mersenne vom 24. Juni 1636
(\cite{00533}\textit{Varia opera mathematica}, Toulouse 1679, S.\;122\textendash124; \textit{OF} II, S.\;17\textendash20).}}
id quod sumsit 
\edtext{Archimedes,\indexN{\textls{Archimedes} 287\textendash212 v. Chr.}}{%
\lemma{Archimedes}%
\Cfootnote{%
\indexN{\textls{Archimedes} 287\textendash212 v. Chr.}\textsc{Archimedes}, 
\cite{01010}\textit{De planorum aequilibris libri duo}, lib.~I, prop.~VI\textendash VII.}}
brachia\indexS{brachia aequalia} aequalia etiam Horizonti\indexS{brachia horizonti non parallela} non parallela in aequilibrio fore pondere existente aequali, non esse rigorose verum; nempe cum directiones\indexS{directiones gravium} gravium non parallelae sunt, sed tendunt ad centrum\indexS{centrum gravium} gravium,
\edtext{\foreignlanguage{french}{\textit{soit le centre\indexS{centre de la terre} de la terre \textit{A}, le grave \textit{E} au point \textit{E}, et le point \textit{N} dans la superficie ou ailleurs plus éloigné du centre que le point \textit{E}, je ne dis pas que le poids \textit{E} pese moins en \textit{E} qu'en \textit{N} mais}}}{%
\lemma{\textit{soit} \bracks{...} \textit{\textit{N} mais}}%
\Cfootnote{%
\cite{02580}\indexN{\textls{Fermat}, Pierre de 1601\textendash1665}\textsc{P. Fermat}, Brief an \indexN{\textls{Mersenne} (Mersennus), Marin 1588\textendash1648}M. Mersenne vom 24. Juni 1636 (\cite{00533}\textit{Varia opera}, S.\;123; \textit{OF} II, S.\;20).}}
\foreignlanguage{french}{que la
\edtext{force en \textit{N}}{%
\lemma{force en}%
\Bfootnote{%
\textit{(1)}~\textit{A}
\textit{(2)}~\textit{N}~\textit{L}}}
par le moyen de la balance le retiendra plus aisement.}
\edtext{\foreignlanguage{french}{\textit{Le principe que M. de Fermat\indexN{\textls{Fermat}, Pierre de 1601\textendash1665} demande}}
\edtext{\foreignlanguage{french}{pour la}}{%
\lemma{pour}%
\Bfootnote{%
\textit{(1)}~\textit{les}
\textit{(2)}~\textit{la}~\textit{L}}}
\foreignlanguage{french}{Geostatique est, que si deux poids\indexS{poids egaux} egaux sont joints par une ligne droite ferme et de soy sans poids, et qu'ainsi disposés, ils puissent descendre librement, ils ne reposeront jamais jusqu'à ce que le milieu de la ligne s'unisse au centre\indexS{centre commun} commun des choses pesantes.}%
}{%
\lemma{\textit{Le principe} \bracks{...} \textit{pesantes}}%
\Cfootnote{%
\indexN{\textls{Pascal}, Étienne 1588\textendash1651}É.\;Pascal und 
\indexN{\textls{Roberval} (Robervallius), Gilles Personne de 1602\textendash1675}G.\;Roberval, 
\cite{02581}Brief an 
\indexN{\textls{Fermat}, Pierre de 1601\textendash1665}P.\;Fermat vom 16. August 1636 
(\cite{00533}\textit{Varia opera}, S.\;124; \textit{OF} II, S.\;35\textendash36; mit Änderungen).}}
\pend
\begin{edfigure}
\includegraphics[scale=1]{%
gesamttex/edit_VIII-4/images/LH_35_14_02_022-023_d1_022r.pdf}
\edcaption{\lbrack\textit{Fig.\;1}\rbrack}
\end{edfigure}
\pstart
Theorema Fermatii\indexN{\textls{Fermat}, Pierre de 1601\textendash1665} tale: \textit{A} centrum\indexS{centrum gravium} gravium seu terrae\indexS{centrum terrae}, erunt \textit{F} et \textit{M} corpora in aequilibrio\indexS{aequilibrium}
\edtext{quando \indexS{arcus}arcus}{%
\lemma{}%
\Bfootnote{%
quando \bars{quando \textit{streicht Hrsg.}} 
arcus~\textit{L}}}
seu distantiae vectis circularis reciproce ut corpora seu 
\edtext{$FH:MH\squaredots M:F$ (+ de reliquo}{%
\lemma{$FH:MH\squaredots M:F$}%
\Bfootnote{%
\textit{(1)}~(+ ipsum tractatum Fermatii\indexN{\textls{Fermat}, Pierre de 1601\textendash1665},
\bars{in \textit{erg.}}
libris iis
\textit{(a)}~adjunctum,
\textit{(b)}~relatum de Geostaticis non reperio
\textit{(2)}~(+ de reliquo~\textit{L}}}
puto pro his quae a centro inaequaliter distant omnia absolvi posse, si consideremus idem esse hic 
\edtext{sive grave sit in \textit{M}, sive in (\textit{M}) et ex \textit{M}}{%
\lemma{sive grave}%
\Bfootnote{%
\textit{(1)}~\textit{M}
\textit{(2)}~sit in \textit{M}, sive in (\textit{M})
\textit{(a)}~unde
\textit{(b)}~et ex \textit{M} ~\textit{L}}}%
pendeat ope 
\edlabel{LH_35_14_02_022r_1a}%
\edtext{}{%
\xxref{LH_35_14_02_022r_1a}{LH_35_14_02_022r_1b}%
\lemma{fili \textit{M}~+).}%
\Bfootnote{%
\href{\Scan{LH_035_14_02_022v}}%
{\lbrack22~v\textsuperscript{o}\rbrack}
\textit{(1)}~(+ Centrum gravitatis corporum non manere semper idem in rigore, sed variari
\textit{(a)}~in si
\textit{(b)}~.~Pro appropinquatione ad centrum terrae, ex eo demonstro
\textit{(2)}~(+ Ponamus~\textit{L}}}%
fili \textit{M}(\textit{M})~+).
\pend
\pstart
\noindent
\href{\Scan{LH_035_14_02_022v}}%
{\lbrack22~v\textsuperscript{o}\rbrack}
\pend
\nobreak
\begin{edfigure}[0em][1em]
\includegraphics[scale=1]{%
gesamttex/edit_VIII-4/images/LH_35_14_02_022-023_d2_022r.pdf}
\edcaption{\lbrack\textit{Fig.\;2}\rbrack}
\end{edfigure}
\pstart
(+ Ponamus\edlabel{LH_35_14_02_022r_1b}
\textit{A} et \textit{B} connecti recta \textit{AB} et esse \textit{C} centrum grav., si jam collocentur in centro terrae, sic ut \textit{C} 
\edtext{coincidat centro terrae. Quaeritur}{%
\lemma{coincidat centro terrae.}%
\Bfootnote{%
\textit{(1)}~Videtur exam
\textit{(2)}~Considera
\textit{(3)}~Quaeritur~\textit{L}}}
an ita maneant et quomodo se collocent. Et examinandum,
an mutatione obtineri possit major vicinia in summa id est sumendum est punctum \textit{C} tale, ut distantiis in corporum quaevis puncta dictis sint summae utrobique aequales. An hic centrum ipsorum globum considerari sufficiat examinandum. Et non videtur, hic enim NB, distantiae non sunt rectae parallelae, ab aliquo axe sed ductae ad ipsum
punctum. NB NB~+).
\pend
\pstart
\edtext{%
\edtext{\textit{Ax. 1.}}{\lemma{}\Bfootnote{\textit{Ax. 1.} \textit{erg.~L}}}
\textit{Si grave\indexS{grave quiescens} quiescens ab aliquo puncto suspendatur, gravitabit super lineam rectam, punctum\indexS{punctum suspensionis} suspensionis et centrum terrae conjungentem, alioqui grave utique non quiesceret.}}{%
\lemma{\textit{Ax. 1.} \bracks{...} \textit{quiesceret}}%
\Cfootnote{%
\indexN{\textls{Fermat}, Pierre de 1601\textendash1665}%
\textsc{P.\;Fermat}, \cite{02582}Brief an \indexN{\textls{Pascal}, Étienne 1588\textendash1651}É.\;Pascal und \indexN{\textls{Roberval} (Robervallius), Gilles Personne de 1602\textendash1675}G.\;Roberval vom 23. August 1636 (\cite{00533}\textit{Varia opera}, S.\;131; \textit{OF} II, S.\;52; mit Änderungen im letzten Teilsatz).}}
\pend
\begin{edfigure}[0em][1em]
\includegraphics[scale=1]{%
gesamttex/edit_VIII-4/images/LH_35_14_02_022-023_d3_022v.pdf}
\edcaption{\lbrack\textit{Fig.\;3}\rbrack}
\end{edfigure}
\pstart
\edtext{\textit{Axiom. 2. In vecte circulari}
\edtext{\textit{cujus} centrum est centrum terrae et}{%
\lemma{}%
\Bfootnote{\textit{cujus} centrum est centrum terrae et \textit{erg.~L}}}
\textit{cujus dimidium sit punctum suspensionis}
\edtext{\textit{N,}}{\lemma{}\Bfootnote{\textit{N} \textit{erg.~L}}}
\textit{si ex utraque in punctis aequalium sectionum} \textit{D} et \textit{G} \textit{gravia aequalia collocentur corpus quiescet}.%
}{%
\lemma{\textit{Axiom. 2.} \bracks{...} \textit{quiescet}}%
\Cfootnote{%
\cite{02582}a.\,a.\,O., mit Auslassungen und Änderungen.}}
\pend
\begin{edfigure}[.5em][1em]
\includegraphics[scale=1]{%
gesamttex/edit_VIII-4/images/LH_35_14_02_022-023_d4_022v.pdf}
\edcaption{\lbrack\textit{Fig.\;4}\rbrack}
\end{edfigure}
\pstart
\edtext{%
\textit{Axiom. 3. Si vero punctum suspensionis inaequaliter vectem dividat non manebit} in aequilibrio \textit{sed inclinabitur ex parte}
\edtext{\textit{majoris portionis.}}{%
\lemma{majoris}%
\Bfootnote{%
\textit{(1)}~portionibus
\textit{(2)}~portionis~\textit{L}}}%
}{%
\lemma{\textit{Axiom. 3.} \bracks{...} \textit{portionis}}%
\Cfootnote{\cite{02582}a.\,a.\,O., mit Auslassungen und Änderungen.}}
\edlabel{LH_35_14_02_022-023_6a}%
\edtext{}{%
\xxref{LH_35_14_02_022-023_6a}{LH_35_14_02_022-023_6b}%
\lemma{Hinc \bracks{...} anguli}%
\Cfootnote{%
\cite{02582}a.\,a.\,O. (\cite{00533}\textit{Varia opera}, S.\;131\,f.; \textit{OF} II, S.\;52\,f.).}}%
Hinc ad modum Archimedeum\indexN{\textls{Archimedes} 287\textendash212 v. Chr.}
\edtext{demonstratur (corpora in vecte aequaliter distribuendo) esse aequilibrium,}{%
\lemma{demonstratur}%
\Bfootnote{%
\bars{(~\textit{gestr. u. wieder gültig gemacht}}
\textit{(1)}~disponendo
\textit{(2)}~corpora in vecte aequaliter distribuendo)
\textit{(a)}~manente
\textit{(b)}~esse~\textit{L}}}
si
\edtext{sit \textit{DE} ad \textit{EG},}{%
\lemma{sit}%
\Bfootnote{%
\textit{(1)}~\textit{GE}
\textit{(2)}~\textit{EG}
\textit{(3)}~\textit{DE} ad \textit{EG}~\textit{L}}}
ut \textit{G} ad \textit{D}. Hinc sequitur corpora aequidistantia\indexS{corpora aequidistantia} a centro terrae aequaliter ponderare si
\edtext{anguli\indexS{angulus radii} radiorum a centris eorum}{%
\lemma{anguli}%
\Bfootnote{%
\textit{(1)}~cum
\textit{(2)}~radiorum 
\textit{(a)}~ab extrem
\textit{(b)}~a centris~\textit{L}}}
ad centrum terrae ductorum faciant cum radio ex puncto suspensionis ad centrum terrae
\edtext{ducto angulos}{%
\lemma{ducto}%
\Bfootnote{%
\textit{(1)}~radios
\textit{(2)}~angulos~\textit{L}}}
corporibus reciproce proportionales seu si corpora duo
\edtext{suspensa ex}{\lemma{}\Bfootnote{suspensa \bars{sint \textit{streicht Hrsg.}} ex~\textit{L}}}
\textit{D} et \textit{G} sint
\edtext{in libra\indexS{libra}}{%
\lemma{in}%
\Bfootnote{%
\textit{(1)}~vecte
\textit{(2)}~libra~\textit{L}}}
\textit{DG}, sitque punctum suspensionis \textit{F}, et ang. \textit{GNF} ad ang. \textit{DNE}, ut \textit{D} ad \textit{G} erunt in
\edtext{aequilibrio. Quod si}{%
\lemma{aequilibrio.}%
\Bfootnote{%
\textit{(1)}~Haec intellige
\textit{(2)}~Quod~\textit{L}}}
punctum \textit{F} cadat in ipsum centrum
\edtext{gravium \textit{N}, tunc evanescent anguli.\edlabel{LH_35_14_02_022-023_6b}}{%
\lemma{gravium \textit{N},}%
\Bfootnote{%
\textit{(1)}~tunc \textit{DN} et \textit{GN} cadent
\textit{(2)}~tunc~\textit{L}}}
\edtext{(+~Hic}{%
\lemma{(+}%
\Cfootnote{Die Klammer bleibt offen.}}
Fermatius\indexN{\textls{Fermat}, Pierre de 1601\textendash1665} rem peculiari ratione tractat, ego vero ex principio priore
\edtext{suppon\bracks{o}}{%
\lemma{}%
\Bfootnote{%
suppone
\textit{L ändert Hrsg.}}}
\textit{FN} esse infinite parvum. Hinc invenio
\edtext{difficultatem, nam deberet}{%
\lemma{difficultatem}%
\Bfootnote{%
\textit{(1)}~deberet
\textit{(2)}~nam deberet~\textit{L}}}
prodire
\edtext{ve\bracks{ra}}{%
\lemma{}%
\Bfootnote{%
veia
\textit{L ändert Hrsg.}}}
regula in casu nostro nempe
\edtext{debere\bracks{t}}{%
\lemma{}%
\Bfootnote{%
debere~\textit{L ändert Hrsg.}}}
esse \textit{FG} (seu \textit{NG})
\edtext{et \textit{FD} ut \textit{D} ad \textit{G} cum tamen faciendo angulos\indexS{anguli reciproce proportionales} reciproce proportionales}{%
\lemma{et \textit{FD}}%
\Bfootnote{%
\textit{(1)}~reciproce proportionalia ipsis
\textit{(2)}~ut \textit{D} ad \textit{G} cum tamen
\textit{(a)}~ductis
\textit{(b)}~faciendo~\textit{L}}}
fiant aequales partes.
\edtext{Unde colligo}{%
\lemma{}%
\Afootnote{%
\textit{Am Rand}: NB}}
pro certo errorem alicubi subesse in argumento Fermatii\indexN{\textls{Fermat}, Pierre de 1601\textendash1665} seu esse. Ego igitur sic ratiocinor.
\pend
\begin{edfigure}
\includegraphics[scale=1]{%
gesamttex/edit_VIII-4/images/LH_35_14_02_022-023_d5_022v.pdf}
\edcaption{\lbrack\textit{Fig.\;5}\rbrack}
\end{edfigure}
\pstart
Ponamus fieri motum in alterutram partem et \textit{G} descendere.
\edtext{Debet omnibus compensatis}{%
\lemma{Debet}%
\Bfootnote{%
\textit{(1)}~plus jam m
\textit{(2)}~omnibus~\textit{L}}}
plus esse descensus\indexS{descensus} quam ascensus\indexS{ascensus}, quod si et aequalis nulla erat ratio et omnia erant in aequilibrio\indexS{aequilibrium}. Sit igitur \textit{ND = NG}
\edtext{primum.}{%
\lemma{primum}%
\Bfootnote{%
\textit{(1)}~, et sit =
\textit{(2)}~erit 
\textit{(a)}~momentum
\textit{(b)}~ipsius \textit{D} distantia a centro ducta in \textit{D}, seu
\textit{(aa)}~momentum 
\textit{(bb)}~quantitas distantiae materiae, \textit{ND} et similiter quantitas distantiae ipsius
\textit{(aaa)}~\textit{g}
\textit{(bbb)}~\textit{G}
\textit{(aaaa)}~sit in \textit{g}
\textit{(bbbb)}~quaeratur differentia inter \textit{ACG}
\textit{(3)}~. Ponatur~\textit{L}}}
Ponatur fieri descensus ex \textit{G} in (\textit{G}) et Ascensus ex \textit{D} in (\textit{D}). Ajo ut sit aequilibrium, debere tale sumi punctum \textit{F}, ut sint saltem in initio seu infinite parvis ipsae appropinquationes seu differentiae\indexS{differentiae reciproce proportionales} inter \textit{NG} et \textit{N}(\textit{G}) item \textit{ND} et \textit{N}(\textit{D}) ipsis corporibus reciproce proportionales. Investigentur calculo.
\pend
\vspace{0.5em}
\pstart
\noindent
\lbrack\textit{Nachfolgend kleingedruckter Text samt \bracks{\textit{Fig.\;6}} gestrichen:}\rbrack %
\pend
\nobreak
\footnotesize	%
\begin{edfigure}[.5em][.5em]
\includegraphics[scale=1]{%
gesamttex/edit_VIII-4/images/LH_35_14_02_022-023_d6_022v.pdf}
\edcaption{\normalsize\lbrack\textit{Fig.\;6, gestrichen}\rbrack}
\end{edfigure}
\pstart 
\noindent
Sit $ND=t$ et $NG=e$ (abstrahamus nunc ab aequalitate\bracks{)} et \textit{NR} data
\edtext{\textit{r} et \textit{FR} data}{%
\lemma{\textit{r} et}%
\Bfootnote{%
\textit{(1)}~\textit{fr} d
\textit{(2)}~\textit{FR} data~\textit{L}}}
\textit{f} et \textit{RD\;\greek{d}}, et \textit{RG\;g} erit
\edtext{$tt=\delta\delta+rr$, et $ee=gg+rr$ et $FD=\delta+f$}{%
\lemma{$tt=\delta\delta+rr$, et}%
\Bfootnote{%
\textit{(1)}~\textit{NG}
\textit{(2)}~\textit{NE}
\textit{(3)}~$ee=gg+rr$~\textit{L}
\textit{(a)}~fiet
\textit{(b)}~\bars{et \textit{streicht Hrsg.}} $FD=\delta+f$~\textit{L}}}
\edtext{et $FG=g-f$}{%
\lemma{}%
\Bfootnote{et $FG=g-f$ \textit{erg.~L}}}
et $DB=b$ et $BR=\sqrt{FD^2-bb}=(D)Q$. et $N(D)=\surd\llcorner rr+2rb\ovalbox{+bb}+FD^2\ovalbox{-bb}\lrcorner$ seu $ND=\sqrt{rr+2rb+\delta^2+2\delta f+f^2}$ et $N(D)-ND=\sqrt{rr+2rb+\delta^2+2\delta f+f^2}-\sqrt{\delta\delta+rr}$. Ob infinite parva autem
\edtext{sufficiet quaeri}{%
\lemma{}%
\Bfootnote{%
sufficiet \bars{ad \textit{streicht Hrsg.}} 
quaeri~\textit{L}}}
differentialem ipsius $\sqrt{\delta\delta+rr}$, quae est
\edlabel{LH_35_14_02_022-023_5a}%
\edtext{}{%
\xxref{LH_35_14_02_022-023_5a}{LH_35_14_02_022-023_5b}%
\lemma{est $dt=$}%
\Bfootnote{%
\textit{(1)}~$2\delta\, d\delta$
\textit{(2)}~$\cancel{2}\delta\,d\overline{\delta}:\cancel{2}\sqrt{\delta\delta+rr}$
\textit{(a)}~ergo
\textit{(b)}~diff \textit{ND}
\textit{(c)}~$d\delta=FD-FB$
\textit{(d)}~Similiter $de=$
\textit{(aa)}~$g\,\delta g$
\textit{(bb)}~$g\,dg:\sqrt{gg+rr}$. 
\textit{(aaa)}~Ubi patet
\textit{(bbb)}~Debent~\textit{L}}}%
$dt=\cancel{2}\delta\,d\overline{\delta}:\cancel{2}\sqrt{\delta\delta+rr}$.
\pend
\pstart
Similiter $de=g\,dg:\sqrt{gg+rr}$. Debent\edlabel{LH_35_14_02_022-023_5b}
autem
\edtext{esse $d\overline{e}$ in pond. $G=d\overline{\delta}$}{%
\lemma{esse}%
\Bfootnote{%
\textit{(1)}~$d\overline{e}\, G=d\overline{\delta}\, D$
\textit{(2)}~$d\overline{e}$~\textit{L}}}
in
\edtext{pondus \textit{D} et cum}{%
\lemma{pondus \textit{D}}%
\Bfootnote{%
\textit{(1)}~. Ergo fiet
\textit{(2)}~et cum~\textit{L}}}
$d\overline{\delta}$ et $d\overline{g}$ seu \textit{DB} et \textit{GH} sint sinus\indexS{sinus anguli} versi eorundem angulorum \textit{DF}(\textit{D}) et \textit{GF}(\textit{G}) erunt ut radii, seu ut \textit{FD} ad \textit{FG}. Ergo $d\overline{\delta}:d\overline{g}\squaredots FD$ seu $\delta +f:FG\;\overline{g-f}$ seu $d\overline{\delta}=d\overline{g}\;\overline{\delta +f}:\overline{g-f}$. Ergo fiet
\edtext{tandem $D\cdot\overline{\delta +f}:\overline{\sqrt{\delta\delta+rr}\cdot\overline{g-f}}=G\cdot\sqrt{gg+rr}$}{%
\lemma{tandem $D\,\cdot$}%
\Bfootnote{%
\textit{(1)}~$\delta\, d\overline{g}$
\textit{(2)}~$\overline{\delta +f}$
\textit{(a)}~$h\,\overline{g-f}=$
\textit{(b)}~$=g\, d\overline{g}\,\delta\delta$
\textit{(c)}~$:\overline{\sqrt{\delta\delta+rr}\cdot\overline{g-f}}=G\cdot\sqrt{gg+rr}$.~\textit{L}}}
seu $D:G=\overline{g-f}\cdot\sqrt{\delta\delta+rr}:\overline{\overline{\delta +f}\cdot\sqrt{gg+rr}}$.
Hinc si \textit{r} sit infinita fit $D:G\squaredots \overline{g-f}:\overline{g+f}$ seu $\squaredots\,FG:DG$. ut dudum notum, et sic consentit. Hoc Fermatius\indexN{\textls{Fermat}, Pierre de 1601\textendash1665} ex suo ducere non posset. Hinc theorema
\edtext{in aequilibrio:}{%
\lemma{}%
\Bfootnote{in aequilibrio \textit{erg.~L}}}
\textit{D} ad \textit{G} in composita ratione ex directa distantiarum a centro gravium, et reciproca distantiarum a centro aequilibrii. Hinc in casu vectis in centro posito vires sunt ut quantitates cessantibus momentis.
\pend
\normalsize		%
\vspace{0.5em}
\pstart
\href{\Scan{LH_035_14_02_023r}}%
{\lbrack23~r\textsuperscript{o}\rbrack}
\edtext{Theorema meum}{%
\lemma{}%
\Bfootnote{%
Theorema \bars{igitur \textit{gestr.}} 
meum~\textit{L}}%
\lemma{}%
\Afootnote{%
\textit{Am Rand:} Verum Theorema sed malus fuit calculus quem pag. seq. correxi.}}
\edtext{est: Duo}{%
\lemma{est:}%
\Bfootnote{%
\textit{(1)}~in aequilibrio 
\textit{(2)}~Duobus
\textit{(a)}~grav
\textit{(b)}~punctis gravibus
\textit{(3)}~Duo~\textit{L}}}
puncta gravia in alicujus vectis rectilinei extremitate posita aequilibrium faciunt, si sint eorum gravitates in ratione\indexS{ratio composita} composita
\edtext{ex ratione}{%
\lemma{ex}%
\Bfootnote{%
\textit{(1)}~distant
\textit{(2)}~ratione~\textit{L}}}
distantiarum a centro gravium directa,
\edtext{et distantiarum\indexS{distantia a centro suspensionis}}{%
\lemma{et}%
\Bfootnote{%
\textit{(1)}~centro
\textit{(2)}~aequi
\textit{(3)}~ratione
\textit{(4)}~distantiarum~\textit{L}}}
a centro suspensionis
\edtext{reciproca. Hinc}{%
\lemma{reciproca.}%
\Bfootnote{%
\textit{(1)}~Et cor
\textit{(2)}~Hinc~\textit{L}}}
corollaria consequuntur quorum veritas\indexS{veritas}
aliunde comperta est.
\edtext{Nimirum in casu}{%
\lemma{Nimirum}%
\Bfootnote{%
\textit{(1)}~in casu
\textit{(a)}~aequilibr
\textit{(b)}~pa
\textit{(2)}~in casu~\textit{L}}}
parallelismi\indexS{parallelismus} linearum directionis seu communis staticae, cum distantiae\indexS{distantia infinita} a centro gravium sint infinitae,
\edtext{adeoque}{%
\lemma{}%
\Bfootnote{adeoque \textit{erg.~L}}}
et inter se aequales (hoc est differentiam habentes differentibus incomparabilem). Restat sola ratio distantiarum a centro suspensionis reciproca. Idem est in casu aequidistantiae\indexS{aequidistantia} a centro gravium. In casu
\edtext{quo centrum}{%
\lemma{quo}%
\Bfootnote{%
\textit{(1)}~punctum
\textit{(2)}~centrum~\textit{L}}}
\edtext{suspensionis collocatur}{%
\lemma{suspensionis}%
\Bfootnote{%
\textit{(1)}~incidit in punctum
\textit{(2)}~collocatur~\textit{L}}}
in centro gravium oportet gravia
\edtext{esse pondera}{%
\lemma{esse}%
\Bfootnote{%
\textit{(1)}~magnitudine
\textit{(2)}~pondera~\textit{L}}}
aequalia, nec momenta\indexS{momentum} quicquam
\edtext{faciunt. Imo}{%
\lemma{faciunt}%
\Bfootnote{%
\textit{(1)}~, et
\textit{(2)}~. Imo~\textit{L}}}
revera puncti suspensionis nulla amplius ratio habetur; et quantum accedit tantum repellitur.
\pend
\begin{edfigure}[.25em][1em]
\includegraphics[scale=1]{%
gesamttex/edit_VIII-4/images/LH_35_14_02_022-023_d7_022v.pdf}
\edcaption{\lbrack\textit{Fig.\;7}\rbrack}
\end{edfigure}
\pstart
Est tamen in hoc ex calculo derivando subtilis quaedam difficultas, nam si centrum suspensionis fingamus extra centrum gravium, poni res non procedit. An ergo dicemus, cum nulla sit ratio cur alibi, potius hic supponi debere. Hinc etiam patet Axioma aliquod
\edtext{Fermatii\indexN{\textls{Fermat}, Pierre de 1601\textendash1665}}{%
\lemma{Fermatii}%
\Cfootnote{%
Siehe \indexN{\textls{Fermat}, Pierre de 1601\textendash1665}%
\textsc{P.~Fermat},
\cite{02585}\textit{Propositio geostatica}, Beilage zu einem verlorenen Brief an \indexN{\textls{Carcavy} (Carcavi), Pierre de um 1600\textendash1684}Carcavy vom Mai 1636 (\cite{00533}\textit{Varia Opera}, S.\;143\textendash144; \textit{OF} II, S.\;6\textendash11).}}
non esse verum. Gravia\indexS{gravia aequalia} duo aequalia \textit{D}
\edtext{et \textit{G} in extremo}{%
\lemma{et \textit{G}}%
\Bfootnote{%
\textit{(1)}~non
\textit{(2)}~in extremo~\textit{L}}}
rectae \textit{DG} non gravis non quiescere nisi medium eorum incidat in centrum gravium.
\edtext{Nam modo}{%
\lemma{Nam}%
\Bfootnote{%
\textit{(1)}~etsi
\textit{(2)}~modo~\textit{L}}}
centrum gravium in rectam ipsam \textit{DG} cadat, etiamsi alterutrum in centro ponatur tamen quiescent.
\pend
\pstart
Sed ubi a punctis
\edtext{in recta per centrum tendente positis}{%
\lemma{}%
\Bfootnote{in recta per centrum tendente positis \textit{erg.~L}}}
abitur, jam conatus\indexS{conatus} ad centrum gravium tendendi non sunt ut gravia, nec corporum vires\indexS{vis corporis} ut magnitudines\indexS{magnitudo corporum}. Sed puncti cujusque momentum ex generali theoremate quaeri debet. Et primum quaerendum momentum
\edtext{rectae\indexS{momentum rectae}, et quadrati\indexS{momentum quadrati}}{%
\lemma{rectae,}%
\Bfootnote{%
\textit{(1)}~et plani, et cubi
\textit{(2)}~et quadrati~\textit{L}}}
et cubi\indexS{momentum cubi}, quia his complentur aliae figurae.
Quanquam nec hoc necesse sit, \edtext{quia tunc}{%
\lemma{quia}%
\Bfootnote{%
\textit{(1)}~sunt
\textit{(2)}~tunc~\textit{L}}}
tantum rectis quadratis planis infinite parvis 
\edlabel{LH_35_14_02_022-023_4a}%
\edtext{}{%
\xxref{LH_35_14_02_022-023_4a}{LH_35_14_02_022-023_4b}%
\lemma{opus.}%
\Bfootnote{%
\textit{(1)}~Fermatius\indexN{\textls{Fermat}, Pierre de 1601\textendash1665} putat esse pondera reciproce ut angulos modo explicato
\textit{(2)}~Sequitur Robervallii
\textit{(a)}~ratiocinatio quam
\textit{(b)}~Epistola~\textit{L}}}%
opus.
\pend
\pstart
Sequitur
\edtext{Robervallii\indexN{\textls{Roberval} (Robervallius), Gilles Personne de 1602\textendash1675} Epistola\edlabel{LH_35_14_02_022-023_4b}}{%
\lemma{Robervallii Epistola}%
\Cfootnote{%
\indexN{\textls{Roberval} (Robervallius), Gilles Personne de 1602\textendash1675}%
\cite{02583}\textsc{G. Roberval}, Brief an \indexN{\textls{Fermat}, Pierre de 1601\textendash1665}P. Fermat vom 11. Oktober 1636 (\cite{00533}\textit{Varia opera}, 
S.\;138\textendash141; \textit{OF} II, S.\;75\textendash83).}}
et ratiocinatio quam examinemus; 
\edlabel{LH_35_14_02_022-023_3a}%
\edtext{}{%
\xxref{LH_35_14_02_022-023_3a}{LH_35_14_02_022-023_3b}%
\lemma{sit}%
\Bfootnote{%
\textit{(1)}~Quia
\textit{(2)}~Si~\textit{L}}}%
sit\bracks{:}
\pend
\begin{edfigure}
\includegraphics[scale=1]{%
gesamttex/edit_VIII-4/images/LH_35_14_02_022-023_d8_023r.pdf}
\edcaption{\lbrack\textit{Fig.\;8}\rbrack}
\end{edfigure}
\pstart
Si\edlabel{LH_35_14_02_022-023_3b}
\edtext{potentiae \textit{B} et \textit{C} aequales, et anguli directionum ad vectem \textit{ECA}, \textit{DBA} aequales erit vis\indexS{vis aequalis} aequalis et in aequilibrio seu opposita et in \textit{B} et \textit{F} dupla.}{%
\lemma{Si potentiae \bracks{...} dupla}%
\Cfootnote{\cite{02583}a.\,a.\,O. (\cite{00533}\textit{Varia opera}, S.\;138; \textit{OF} II, S.\;75; Französisch mit Änderungen).}}
Si brachia inaequalia et anguli aequales, ut et potentiae effectus erit inaequalis et pro majori brachio major. Supponit praeterea autor
\edtext{noster, si eadem sit}{%
\lemma{noster,}%
\Bfootnote{%
\textit{(1)}~si sit
\textit{(2)}~si eadem sit~\textit{L}}}
potentia linea directionis, et centrum\indexS{libra}\indexS{centrum librae} librae nihil referre quod sit brachium\indexS{brachium librae} seu per \textit{B} eodem modo trahi, sive in \textit{F} sive in \textit{C} sive in \textit{G}. Hoc principium ait Robervallius\indexN{\textls{Roberval} (Robervallius), Gilles Personne de 1602\textendash1675} tacite ab autoribus assumi et ab experientia confirmari. Hinc iam Robervallius\indexN{\textls{Roberval} (Robervallius), Gilles Personne de 1602\textendash1675} probat generaliter. Si potentiae\indexS{potentiae contrariae}
\edtext{contrariae}{%
\lemma{}%
\Bfootnote{contrariae \textit{erg.~L}}}
sint reciprocae perpendicularibus ex centro in lineas directionum ductis, fore aequilibrium.
\edtext{\indexN{\textls{Fermat}, Pierre de 1601\textendash1665}Fermatius (fol.\;147) negat}{%
\lemma{Fermatius (fol.\;147) negat}%
\Cfootnote{Siehe \indexN{\textls{Fermat}, Pierre de 1601\textendash1665}\textsc{P.\;Fermat}, \cite{02584}Brief an \indexN{\textls{Roberval} (Robervallius), Gilles Personne de 1602\textendash1675}G.\;Roberval vom 7. Dezember 1636 (\cite{00533}\textit{Varia opera}, S.\;147; \textit{OF} II, S.\;89).}}
hoc principium Robervallii\bracks{,}%
\indexN{\textls{Roberval} (Robervallius), Gilles Personne de 1602\textendash1675}
quod potentia eodem modo
\edtext{trahat in \textit{G} vel \textit{C} vel \textit{F}}{%
\lemma{trahat}%
\Bfootnote{%
\textit{(1)}~per \textit{A} \textit{B} et \textit{C}
\textit{(2)}~in \textit{G} vel \textit{C} vel \textit{F}~\textit{L}}}
esse verum. Dignissima res examinari.
\pend
\begin{edfigure}
\hfill
\includegraphics[scale=1]{%
gesamttex/edit_VIII-4/images/LH_35_14_02_022-023_d9_023r.pdf}
\hfill%
\includegraphics[scale=1]{%
gesamttex/edit_VIII-4/images/LH_35_14_02_022-023_d10_023r.pdf}
\hfill
\edcaption{\lbrack\textit{Fig.\;9}\rbrack 
\hspace*{40mm}%
\edtext{\lbrack\textit{Fig.\;10}\rbrack}{%
\lemma{}%
\Afootnote[nonum]{%
\textit{Unter dem Diagramm} \lbrack \textit{Fig.\;10}\rbrack\textit{, gestrichen}:
\textit{DF} \textit{FG} et \textit{GN} \textit{ND}}}
\hspace*{10mm}%
}%
\end{edfigure}
\vspace{0.5em}
\pstart
\noindent
\lbrack\textit{Nachfolgend kleingedruckter Text samt \bracks{\textit{Fig.\;11}} gestrichen:}\rbrack %
\pend
\nobreak
\footnotesize	%
\begin{edfigure}[.5em][.5em]
\includegraphics[scale=1]{%
gesamttex/edit_VIII-4/images/LH_35_14_02_022-023_d11_023r.pdf}
\edcaption{%
\edtext{\normalsize{\lbrack\textit{Fig.\;11, gestrichen}\rbrack}}{%
\lemma{}%
\Afootnote[nonum]{%
\textit{Über das Diagramm} \lbrack \textit{Fig.\;11}\rbrack\textit{:} Reperi veram et meae coincidentem}}}
\end{edfigure}
\pstart 
Fortasse res vera \edtext{erit potentiae}{%
\lemma{erit}%
\Bfootnote{%
\textit{(1)}~si gravis inter descen
\textit{(2)}~potentiae~\textit{L}}}
eadem semper manet linea directionis ut axes\indexS{axis} fixi \textit{L}
\edtext{et \textit{M}, circa}{%
\lemma{et \textit{M}}%
\Bfootnote{%
\textit{(1)}~per
\textit{(2)}~circa~\textit{L}}}
sua centra acti funibus trahant puncta \textit{B} et \textit{C}, et forte secus in motu libero, grave enim ad centrum aliquod tendens ad \textit{C} affixum, gyrando circa \textit{A}, non manet in recta \textit{CM}.
\pend
\normalsize		%
\vspace{0.5em}
\pstart
\noindent
{\normalsize\href{\Scan{LH_035_14_02_023v}}%
{\lbrack23~v\textsuperscript{o}\rbrack}}%
\pend
\nobreak
\begin{edfigure}[0em][1em]
\includegraphics[scale=1]{%
gesamttex/edit_VIII-4/images/LH_35_14_02_022-023_d12_023v.pdf}
\edcaption{\lbrack\textit{Fig.\;12}\rbrack}
\end{edfigure}
\pstart
Puncta\indexS{punctum grave} gravia \textit{A}
\edtext{et \textit{B} annexa extremis}{%
\lemma{et \textit{B}}%
\Bfootnote{%
\textit{(1)}~in
\textit{(2)}~annexa
\textit{(a)}~extremo
\textit{(b)}~extremis~\textit{L}}}
rectae rigidae pondere carentis \textit{AB} tendunt ad punctum \textit{H}. Quaeritur punctum suspensionis \textit{C}, ut sint in
\edtext{aequilibrio. Ex puncto}{%
\lemma{aequilibrio}%
\Bfootnote{%
\textit{(1)}~ex puncto communi \textit{H}
\textit{(2)}~. Ex centro gra
\textit{(3)}~. Ex puncto~\textit{L}}}
ad quod tendunt \textit{H} agatur perpendicularis in vectem \textit{AB},
\edtext{nempe \textit{HG}, jam \textit{AG} sit \textit{e} et \textit{BG} sit \textit{f}, et}{%
\lemma{nempe}%
\Bfootnote{%
\textit{(1)}~\textit{DF}
\textit{(2)}~\textit{DE}
\textit{(3)}~\textit{HG}
\textit{(4)}~\textit{HG},
\textit{(a)}~vocetur \textit{g}, et \textit{AC} hoc vocetur \textit{h},
\textit{(b)}~jam \textit{AG} sit
\textit{(aa)}~\textit{g}
\textit{(bb)}~\textit{e} et \textit{BG} sit
\textit{(aaa)}~\textit{h}
\textit{(bbb)}~\textit{f}, et~\textit{L}}}
\textit{HG} sit \textit{h}, et \textit{GC}
\edtext{sit \textit{c}, et ex punctis}{%
\lemma{}%
\Bfootnote{%
sit \textit{c}, \bars{fiet \textit{HA}, $\sqrt{rr+gg}=a$ et $H(A)=\sqrt{rr+(g)^2}$ posito $g=(g)=GC$ \textit{gestr.}} 
et ex punctis~\textit{L}}}
(\textit{A}) et (\textit{B}) demittantur in vectem priorem normales\indexS{normalis} (\textit{A})\textit{E}, et (\textit{B})\textit{F}. Compleantur etiam rectangula (\textit{A})\textit{EG}(\textit{G}) et 
\edtext{\bracks{(\textit{B})\textit{FG}((\textit{G}))}.}{%
\lemma{}%
\Bfootnote{%
(\textit{B})\textit{FG}(\textit{G}) 
\textit{L ändert Hrsg.}}}
Jam \textit{HA} vel \textit{H}(\textit{A})
\edtext{vocemus \textit{a}, et \textit{GA}}{%
\lemma{vocemus \textit{a},}%
\Bfootnote{%
\textit{(1)}~et \textit{HB} vel \textit{H}(\textit{B}), \textit{b} et \textit{v}
\textit{(2)}~et
\textit{(a)}~\textit{GH}
\textit{(b)}~\textit{GA}~\textit{L}}}
vel \textit{GE} vocemus \textit{e}, et \textit{HG}
\edtext{vel \textit{H}(\textit{G}) \textit{h}. Similiter \textit{HB} vel \textit{H}(\textit{B}) sit \textit{v}, et \textit{GB} vel \textit{GF} sit \textit{f}, et}{%
\lemma{vel \textit{H}(\textit{G}) \textit{h}.}%
\Bfootnote{%
\textit{(1)}~Similiter \textit{HB} vel \textit{H}(\textit{B}) sit \textit{v} et \textit{HG} vel \textit{A}((\textit{G}))
\textit{(2)}~Similiter \textit{HB} vel \textit{H}(\textit{B}) sit \textit{v}, et
\textit{(a)}~\textit{H}
\textit{(b)}~\textit{GB} vel \textit{GF} sit 
\textit{(aa)}~\textit{H}
\textit{(bb)}~\textit{f}, et~\textit{L}}}
\textit{HG} vel \textit{H}(\textit{GG})
\edtext{sit \textit{H}.}{%
\lemma{}%
\Afootnote{%
\textit{Am Rand:} Mira et notabilis ratio rem ad meam methodum\indexS{methodus tangentium} tangentium revocandi. Eadem quantitas hic diversarum serierum, et ita $H=h$, et tamen $d\overline{H}$ non est aeq. $d\overline{h}$.}}
Fiet $a = \sqrt{ee + hh}$ et $d\overline{a} = \overline{2e\,de + 2h\,dh},\ :2\sqrt{ee + hh}$. Sed $d\overline{e}$ seu \textit{AE} est incomparabilis ipsi \textit{dh}
\edtext{seu (\textit{A})\textit{E}}{%
\lemma{}%
\Bfootnote{seu (\textit{A})\textit{E} \textit{erg.~L}}}
sinus versus recto anguli incomparabilis ad rectum. Ergo fiet $d\overline{a} = h\,dh \sqrt{ee + hh}$. Similiter fiet $v = \sqrt{ff + H^2}$ seu fiet $dv = \overline{2f\,df + 2H\,d\overline{H}} : 2\sqrt{ff + H^2}$. Ubi cum eodem modo elidatur $d\overline{f}$ seu \textit{BF} incomparabilis ipsi (\textit{B})\textit{F} 
\edtext{seu $d\overline{H}$}{%
\lemma{}%
\Bfootnote{seu $d\overline{H}$ \textit{erg.~L}}}
fiet $d\overline{v} = H\,d\overline{H} : \sqrt{ff + H^2}$. Jam pondus in \textit{A} sit \textit{l} et pondus in \textit{B} sit \textit{m}, debet
\edtext{esse $l\,d\overline{a} = m\,d\overline{v}$ seu}{%
\lemma{esse}%
\Bfootnote{%
\textit{(1)}~$mdv=$
\textit{(2)}~$l\,d\overline{a} = m\,d\overline{v}$
\textit{(a)}~rursus
\textit{(b)}~seu~\textit{L}}}
 fiet $lh\,dh\sqrt{ff + H^2} = mH\,dH\sqrt{ee + hh}$. Est autem $H = h$ et \textit{dh} ad \textit{dH} seu (\textit{A})\textit{E} ad (\textit{B})\textit{F}, ut \textit{C}(\textit{A}) ad \textit{C}(\textit{B}) seu ut $\overline{g + e}$ ad $\overline{b - c}$. Fiet $dH = dh\,\overline{f - c}:\overline{e + c}$ adeoque tandem fiet $\cancel{dh}l\cancel{h}\,\overline{e + c}\,\sqrt{ff + hh}=\cancel{dh}m\cancel{h}\,\overline{f - c}\,\sqrt{ee + hh}$ 
\edtext{erit \textit{l} pondus in \textit{A}, seu}{%
\lemma{$\sqrt{ee + hh}$}%
\Bfootnote{%
\textit{(1)}~seu
\textit{(a)}~\textit{l} pon
\textit{(b)}~ut fiat aequilibrium
\textit{(2)}~erit \textit{l} pondus in \textit{A},
\textit{(a)}~ad \textit{m} pondus in \textit{B}, 
\textit{(b)}~ut
\textit{(c)}~in ratione directa
\textit{(d)}~seu~\textit{L}}}
\edtext{$l : m\squaredots\overline{f - c}\,\sqrt{ee + hh}:\overline{e + c}\,\sqrt{ff + hh}$}{%
\lemma{}%
\Afootnote{%
\textit{Am Rand:} calc.}}
seu $l:m \squaredots \overline{CB\cdot HA} : \overline{CA\cdot HB}$, seu
\edtext{\edlabel{LH_35_14_02_022-023_2a}posito aequilibrio}{%
\lemma{}%
\Bfootnote{posito aequilibrio \textit{erg.~L}}}
pondus in \textit{A} est ad pondus in \textit{B}, in
\edtext{ratione composita}{%
\lemma{ratione}%
\Bfootnote{%
\textit{(1)}~directa
\textit{(2)}~composita~\textit{L}}}
ex directa distantiarum\indexS{distantia a centro gravium} a centro gravium
\edtext{\textit{H}}{%
\lemma{}%
\Bfootnote{\textit{H} \textit{erg.~L}}}
(seu \textit{HA} ad \textit{HB})
\edtext{et reciproca}{%
\lemma{et}%
\Bfootnote{%
\textit{(1)}~composita
\textit{(2)}~reciproca~\textit{L}}}
distantiarum a centro aequilibrii \textit{C}, (seu \textit{CB} ad \textit{CA}). Jam ex \textit{C} ducantur \textit{CN}
\edtext{et \textit{CP} normales\indexS{normalis}\edlabel{LH_35_14_02_022-023_2b}}{%
\lemma{et \textit{CP}}%
\Bfootnote{%
\textit{(1)}~perpendiculares
\textit{(2)}~normales~\textit{L}}}
ad lineas directionum \textit{HA} et \textit{HP}. Item ex \textit{G} ad easdem normales\indexS{normalis} \textit{GQ}, \textit{GR}. Primum $CN: GQ \squaredots CA : GA$, et $GQ : GA \squaredots HG : HA$. Ergo $CN : CA \squaredots HG : HA$, quod et initio patebat. Et
\edtext{similiter $CP : CB \squaredots HG : HB$.}{%
\lemma{similiter}%
\Bfootnote{%
\textit{(1)}~$CR:CB$
\textit{(2)}~$CP:CR$
\textit{(3)}~$CP : CB \squaredots HG : HB$.~\textit{L}}}
Ergo
\edtext{tandem $CP : CN \squaredots \overline{CB\cdot HG : HB} : \overline{CA\cdot HG : HA}$}{%
\lemma{tandem}%
\Bfootnote{%
\textit{(1)}~$CN:CB$
\textit{(2)}~\textit{CP}
\textit{(3)}~$CP : CN \squaredots \overline{CB\cdot HG : HB} : \overline{CA\cdot HG : HA}$~\textit{L}}}
seu $CP : CN \squaredots CB\cdot HA : CA\cdot HB$. Ergo $l : m \squaredots CP : CN$, seu
\begin{tabular}{c}
potentiae\\ pondera
\end{tabular}
in casu aequilibrii sunt reciproce ut perpendiculares ex centro aequilibrii ad lineas directionum. Nimirum ut obiter
\edtext{dicam in triangulo}{%
\lemma{dicam}%
\Bfootnote{%
\textit{(1)}~circuli
\textit{(2)}~in triangulo~\textit{L}}}
ex puncto basin dividente ductae ad latera normales\indexS{normalis} sunt in ratione composita directa segmentorum\indexS{ratio segmentorum}, et reciproca laterum\indexS{ratio laterum}. Nullum corpus praeter sphaeram\indexS{sphaera}, in quo grave extra centrum terrae suspensum servet eam quam in principio habuit positionem, ita Fermatius\indexN{\textls{Fermat}, Pierre de 1601\textendash1665}. Eodem judicante errat
\edtext{Ubaldus\indexN{\textls{Del Monte}, Guidobaldo 1545\textendash1607}}{%
\lemma{Ubaldus}%
\Cfootnote{%
\indexN{\textls{Del Monte}, Guidobaldo 1545\textendash1607}%
\textsc{Guidobaldo del Monte},
\cite{01009}\textit{Mechanicorum Liber}, Pesaro 1577, lib. I, prop. III.}}
putans librae\indexS{brachia librae} brachia, etsi non sint parallela horizonti, aequilibrium tamen servare. Ex
\edtext{sententia}{%
\lemma{sententia \indexN{\textls{Fermat}, Pierre de 1601\textendash1665}Fermatii}%
\Cfootnote{%
Siehe
\indexN{\textls{Fermat}, Pierre de 1601\textendash1665}%
\textsc{P.~Fermat},
\cite{02586}\textit{Nova in Mechanicis Theoremata}, Beilage zum Brief an 
\indexN{\textls{Mersenne} (Mersennus), Marin 1588\textendash1648}Mersenne vom 24. Juni 1636 
(\textit{Varia opera}, S.\;142\textendash143; \textit{OF} II, S.\;23\textendash26).}}
Fermatii\indexN{\textls{Fermat}, Pierre de 1601\textendash1665}, quo propius
\edtext{corpus accedit}{%
\lemma{}%
\Bfootnote{%
corpus \bars{corpus \textit{gestr.}} 
accedit~\textit{L}}}
ad centrum\indexS{centrum terrae} terrae, hoc facilius
\edtext{tollitur. 
\edtext{(+~Hoc}{%
\lemma{(+}%
\Cfootnote{Die Klammer bleibt offen.}}%
}{%
\lemma{tollitur.}%
\Bfootnote{%
\textit{(1)}~Posita scilicet
\textit{(2)}~(+~Hoc~\textit{L}}}
etiam ex meis sequitur, posita scil. eadem ratione\indexS{ratio brachiorum librae} brachiorum librae, et distantia unius ponderis a centro, facilitates quo attollitur alterum sunt ut viciniae ejus respectu centri seu distantiis reciproce
\edtext{proportionales, vires}{%
\lemma{proportionales,}%
\Bfootnote{%
\textit{(1)}~resistentiae
\textit{(2)}~vires~\textit{L}}}
enim ponderum\indexS{vis ponderis} ut distantiae.
\pend
\begin{edfigure}[.5em][1em]
\includegraphics[scale=1]{%
gesamttex/edit_VIII-4/images/LH_35_14_02_022-023_d13_023v.pdf}
\edcaption{\edtext{\lbrack\textit{Fig.\;13}\rbrack}{%
\lemma{\lbrack\textit{Fig.\;13}\rbrack}%
\Cfootnote[nonum]{Leibniz hat zunächst die Strecken \textit{AH} und \textit{BX} verlängert, die Verlängerungen jedoch anschließend gestrichen. Die versehentlich stehen gebliebene Bezeichnung \textit{K} für ihren Treffpunkt streicht Hrsg.}}}
\end{edfigure}
\pstart
Ut rem generaliter quaeramus, etiam
\edtext{cum non concurrant}{%
\lemma{cum}%
\Bfootnote{%
\textit{(1)}~puncta
\textit{(2)}~non concurrant~\textit{L}}}
omnes directiones, ut si trahant quaevis potentiae, exempli gratia chordae\indexS{chordae} vel funes\indexS{funis} axibus\indexS{axis} \textit{M} et \textit{X} circumvoluti trahant \textit{A} et \textit{B} ubi sciendum est duas directiones quatuor \textit{AH}, (\textit{A})(\textit{H}), \textit{BX} et (\textit{B})(\textit{X}) non coincidere in idem punctum. 
Igitur ex puncto\indexS{punctum concursus} concursus ipsorum contactuum circuli, quem vocemus \textit{H}, quia infinite parvo invicem intervallo distant\lbrack ,\rbrack\ ducatur perpendicularis \textit{HG}, et
\edtext{similiter $X\gamma$, ita}{%
\lemma{similiter $X\gamma$}%
\Bfootnote{%
\textit{(1)}~. Et ut ante \textit{HA}
\textit{(2)}~, ita~\textit{L}}}
\textit{H} et \textit{X} \lbrack ,\rbrack\ \textit{H} et \textit{X} ita manebunt fixa, jam ut ante
\edtext{sit \textit{HA} vel $H(A), \textit{a}$, \textit{GA} vel \textit{GE}, \textit{e} \textit{CG}, \textit{c} \textit{XB} vel \textit{X}(\textit{B}), \textit{v}}{%
\lemma{sit \textit{HA}}%
\Bfootnote{%
\bars{vel \textit{H}(\textit{A}) \textit{erg.}},
\textit{a}
\textit{(1)}~$XB=v$
\textit{(2)}~\textit{GA} \bars{vel \textit{GE} \textit{erg.}}, \textit{e} \textit{CG}, \textit{c} \textit{XB} \bars{vel \textit{X}(\textit{B}) \textit{erg.}}, \textit{v}~\textit{L}}}
$\gamma B$
\edtext{vel $\gamma F$,}{%
\lemma{}%
\Bfootnote{vel $\gamma F$ \textit{erg.~L}}}
\textit{f}, \textit{HG}
\edtext{vel \textit{H}(\textit{G}),}{%
\lemma{}%
\Bfootnote{vel \textit{H}(\textit{G}) \textit{erg.~L}}}
\textit{h},
\edtext{et $\bracks{X}\gamma$ vel \bracks{$X(\gamma)$}, \textit{x}. $aa = ee + hh$\bracks{,}}{%
\lemma{}%
\Bfootnote{%
et
\bars{$H\gamma$ \textit{ändert Hrsg.}}
\bars{vel $H(\gamma)$ \textit{erg., ändert Hrsg.}}
\textit{x}. \bars{\textit{CA} vel \textit{gestr.}} 
$aa = ee + hh$~\textit{L}}}
$d\overline{a} = h\,dh: \sqrt{ee + hh}$. Et $vv = ff + xx$. Et $dv = x\,dx: \sqrt{ff + xx}$. Est 
\edtext{autem \textit{dh} (seu (\textit{A})\textit{E}) ad \textit{dx} (seu (\textit{B})\textit{F}) ut \textit{CA} ad \textit{CB}. Ergo $d\overline{a}$ ad $d\overline{v}$}{%
\lemma{autem}%
\Bfootnote{%
\textit{(1)}~$dh:dx\squaredots CA:CB$
\textit{(2)}~ \textit{dh} \bracks{...} \textit{CB}.
\textit{(a)}~seu ut jam vires sunt \textit{l} in $d\overline{a}$ et \textit{m} in $d\overline{v}$.
\textit{(b)}~Ergo $d\overline{a}$ ad $d\overline{v}$~\textit{L}}}
ut $\overline{HG : HA}$ ad
\edtext{$\overline{X\gamma : XB}\cdot\overline{CB : CA}$ seu}{%
\lemma{}%
\Bfootnote{%
 $\overline{X\gamma : XB}\cdot\overline{CB : CA}$ \bars{seu vires ponderum sunt \textit{gestr.}} 
seu~\textit{L}}}
$d\overline{a}$ ad $d\overline{v}$ ut $\overline{HG\cdot CA : HA}$ ad $\overline{X\gamma\cdot XB : CB}$. Jam ductis \textit{CN} et \textit{CP} normalibus in \textit{HA} et \textit{XB} fient $CN : CA \squaredots HG : HA$ seu $CN = HG\cdot CA : HA$. Et eodem modo
\edtext{$CP = X\gamma\cdot \lbrack CB : XB\rbrack$.}{%
\lemma{}%
\Bfootnote{%
$CP = X\gamma\cdot XB : CB$ 
\textit{L ändert Hrsg.}}}
Ergo $d\overline{a} : dv \squaredots CN : CP$. Seu vires\indexS{vires ponderum} ponderum sunt in composita ratione ipsorum ponderum et normalium\indexS{normalis} in lineas directionis. Eadem demonstratio foret etiamsi \textit{H} et \textit{X} non essent puncta fixa, sed ipsa durante motu mobilia. Saltem enim invenietur punctum concursus\indexS{vires ponderum} \textit{H} duarum directionum in \textit{A} et in (\textit{A}) et eadem erit ratiocinatio quae ante. Sed hic notandum hoc modo corporis integri
\edtext{agentis in extremitate non}{%
\lemma{agentis}%
\Bfootnote{%
\textit{(1)}~non
\textit{(2)}~in extremitate non~\textit{L}}}
esse
\edtext{sumend\lbrack as\rbrack}{%
\lemma{}%
\Bfootnote{%
sumendum
\textit{L ändert Hrsg.}}}
\edtext{gravitat\lbrack e\rbrack s\indexS{gravitas},}{%
\lemma{}%
\Bfootnote{%
gravitatis
\textit{L ändert Hrsg.}}}
nisi ejus omnium punctorum directiones parallelas.
\pend
\vspace{.5em}%
\pstart
\noindent
\lbrack\textit{Nachfolgender Text samt \bracks{\textit{Fig.\;14}} befindet sich am Rand von Bl.\;%
\href{\Scan{LH_035_14_02_023v}}%
{23~v\textsuperscript{o}}, auf der Höhe von \refpassage{LH_35_14_02_022-023_2a}{LH_35_14_02_022-023_2b}:}\rbrack
\pend

\begin{edfigure}[0em][1em]
\includegraphics[scale=1]{%
gesamttex/edit_VIII-4/images/LH_35_14_02_022-023_d14_023v.pdf}
\edcaption{\lbrack\textit{Fig.\;14}\rbrack}
\end{edfigure}
\pstart
\edtext{Dato corpore}{%
\lemma{Dato}%
\Bfootnote{%
\textit{(1)}~puncto
\textit{(2)}~corpore~\textit{L}}}
\textit{A} ad centrum terrae tendente, quaeritur quae sit linea directionis \textit{BC} seu quod punctum sustentationis \textit{C}, v.\,g. si ponamus esse triangulum scalenum ut cujus latera 3.4.5.
\pend
